% TOP_PROBLEM: {"id": 5500013, "title": "Point Location in 3D Subdivision", "category_id": 15, "category": {"id": 15, "name": "computer_science", "display_name": "Computer Science", "description": "Computational complexity, algorithms, and theoretical CS.", "slug": "computer-science", "order_index": 15, "created_at": "2026-07-31T15:26:25.671Z"}, "status": "open", "problem_number": "AMR-054-0013", "set_id": 13}
% FIRST_POSTED: 2026-10-01
% ENTRY_KIND: statement_only
% UnsolvedMath ID 5500013; source catalogue code AMR-054-0013
% !TeX program = lualatex
% Definition and sources only; this file is not an LLM solution attempt.
\documentclass[11pt,a4paper]{article}
\usepackage[margin=25mm]{geometry}
\usepackage{fontspec}
\setmainfont{lmroman10-regular.otf}[BoldFont=lmroman10-bold.otf,
 ItalicFont=lmroman10-italic.otf,BoldItalicFont=lmroman10-bolditalic.otf]
\usepackage{amsmath,amssymb,mathtools}
\usepackage{unicode-math}
\setmathfont{latinmodern-math.otf}
\AtBeginDocument{\let\setminus\smallsetminus}
\usepackage{xurl,hyperref}
\hypersetup{colorlinks=true,urlcolor=blue}
\setcounter{secnumdepth}{0}
\setlength{\parindent}{0pt}
\setlength{\parskip}{5pt}
\setlength{\emergencystretch}{3em}
\begin{document}
\section*{Point Location in 3D Subdivision}\label{5500013}
{\small UnsolvedMath ID 5500013 · AMR-054-0013 · Computer Science · AMR Open Problem Lists}
\subsection{Definitions and mathematical statement}
A polyhedral subdivision of three-dimensional space is a finite
face-to-face complex whose three-dimensional cells cover a prescribed
polyhedral domain (with an exterior cell if needed). Its faces, edges,
vertices, and incidence relations are given explicitly. Let $n$ denote
its total boundary-description complexity; in a triangulated description
this is the number of constant-size boundary faces up to constant factors.

A point-location query supplies $q\in\mathbb R^3$ and asks which cell
contains it. Boundary points may be handled by a fixed deterministic
tie rule, or the query can report an incident cell. The central question
is whether every such subdivision admits a data structure with
\[
                    O(n)\text{ storage},\qquad
                    O(\log n)\text{ worst-case query time}.
\]
Preprocessing must construct the structure from the subdivision; its
cost should be stated separately. The bounds refer to a standard
real-arithmetic/pointer model with constant-size geometric records.

Counting each polygonal face as one unit regardless of an unbounded
number of vertices would conceal input complexity, so that convention
is not used. The requested simultaneous space and query bounds are
stronger than obtaining either bound alone with an unrestricted cost
for the other.
\subsection{Short English statement}
Can arbitrary three-dimensional polyhedral subdivisions support
logarithmic-time point location using only space proportional to their
input complexity?
\subsection{Sources}
\begin{itemize}
\item[S1] ULAM.AI and the UnsolvedMath Contributors, supplied problems.json, record 5500013 (AMR-054-0013), AMR Open Problem Lists. Catalogue identity and source transcription; CC BY 4.0.
\url{https://huggingface.co/datasets/ulamai/UnsolvedMath}.
\item[S2] Open Problems Project - Computational Geometry, Problem 13. Original source indexed by AMR-054-0013.
\url{https://topp.openproblem.net/p13}.
\end{itemize}
\end{document}
